% Generated draft: requires mathematical and source review.
\documentclass[11pt]{article}
\usepackage[margin=25mm]{geometry}
\usepackage{fontspec}
\setmainfont{Latin Modern Roman}
\usepackage{amsmath,amssymb,mathtools,unicode-math}
\setmathfont{Latin Modern Math}
\usepackage{enumitem,needspace,xurl,hyperref}
\setlength{\parindent}{0pt}
\setlength{\parskip}{5pt}
\setlength{\emergencystretch}{3em}
\newcommand{\N}{\mathbb{N}}
\newcommand{\Z}{\mathbb{Z}}
\newcommand{\Q}{\mathbb{Q}}
\newcommand{\R}{\mathbb{R}}
\newcommand{\C}{\mathbb{C}}
\newcommand{\F}{\mathbb{F}}
\newcommand{\A}{\mathbb{A}}
\newcommand{\PP}{\mathbb{P}}
\newcommand{\E}{\mathbb{E}}
\newcommand{\eps}{\varepsilon}
\newcommand{\dd}{\,\mathrm{d}}
\begin{document}
\section*{UnsolvedMath ID: 30000436}
\textbf{Source problem code:} catalogue entry
\subsection{Definitions and mathematical statement}
Let $P\subset\R^3$ be a convex polytope and equip its boundary with the intrinsic Euclidean polyhedral metric. Let $T$ be a prescribed geodesic triangulation of $\partial P$ using the vertices of $P$, with no added Steiner vertices. Its edges are intrinsic geodesic arcs on the surface and may cross edges of the original polytope; they need not be original polytope edges.

In the usual single-net formulation, the question asks whether for every such pair $(P,T)$ there is a spanning cut tree $C\subset T^{(1)}$ such that cutting $\partial P$ along $C$ produces a disk admitting a planar development
\[
 D:(\partial P)_{\mathrm{cut}\,C}\longrightarrow\R^2
\]
that is isometric on each triangle, matches adjacent triangles across uncut edges, and has no overlap of triangle interiors. Boundary contact is allowed in this standard nonoverlap convention. The essential quantifiers are
\[
 \forall P\ \forall T\ \exists C\text{ using only edges of the prescribed }T.
\]
Choosing a new favorable triangulation or cutting along unrelated geodesics would be a different, weaker existence problem.

The source's short wording does not explicitly state the cut-tree, single-connected-net or boundary-contact conventions, nor whether geodesic edges must be globally shortest or only intrinsic local geodesics. The displayed tree development is the standard precise interpretation. Allowing arbitrary disconnected pieces would trivialize the question by separating all triangles and is not the intended unfolding target.

The literature assessment checked on 21 August 2026 says the arbitrary-geodesic-triangulation problem remains unresolved beyond special cases such as the tetrahedron. Those special cases and results allowing other cuts do not establish the displayed universal assertion. This is dated supplied triage requiring independent verification.
\subsection{Short English statement}
Can every prescribed geodesic triangulation of a convex polytope surface be cut along its own edges to make one nonoverlapping planar net? The triangulation is fixed, and its geodesic edges need not be the original polytope edges.
\subsection{Sources}
\begin{itemize}
\item \emph{Discrete Differential Geometry} (2006), Oberwolfach Reports, DOI 10.4171/owr/2006/12, pp. 653--728, \url{https://doi.org/10.4171/owr/2006/12}.
\item Literature assessment checked 21 August 2026: original OWR problem, \url{https://doi.org/10.4171/owr/2006/12}; ``convex polytope unfolding status,'' \url{https://jeffe.web.engr.illinois.edu/open/unfold.html}.
\item Local UnsolvedMath registry, permanent ID $30000436$, exact source code \texttt{catalogue entry}, created 25 August 2026 and updated 16 September 2026. The supplied statement is labelled corrected and verified against the official report with source-level content unchanged.
\end{itemize}
\end{document}
